% ============================================================
%  Section 2: Mathematical Modeling of the Non-linear CSTR and Controllability Limits
% ============================================================

\subsection{Non-linear Physical Plant Model}
A continuous stirred tank reactor (CSTR) conducting an irreversible, liquid-phase exothermic reaction $A \xrightarrow{k} B$ is considered \cite{luyben2007chemical, seborg2010process}. Rather than adopting the frequent simplification of instantaneous jacket dynamics, the cooling jacket is modeled explicitly as a dynamic lumped thermal system coupled through non-linear heat convection across the reactor wall \cite{uppal1974dynamic}.

The dynamic behavior of the process is governed by three coupled non-linear ordinary differential equations (ODEs) derived from dynamic species mass and thermal energy balances:
\begin{align}
    \frac{dC_A}{dt} &= \frac{q}{V}(C_{Af} - C_A) - r_A(C_A, T) \label{eq:cstr_ca} \\
    \frac{dT}{dt}   &= \frac{q}{V}(T_f - T) + \frac{(-\Delta H)}{\rho C_p} r_A(C_A, T) - \frac{UA}{V \rho C_p}(T - T_j) \label{eq:cstr_t} \\
    \frac{dT_j}{dt} &= \frac{q_j}{V_j}(T_{jf} - T_j) + \frac{UA}{V_j \rho_j C_{pj}}(T - T_j) \label{eq:cstr_tj}
\end{align}
where $C_A$ denotes the concentration of reactant $A$ in the reactor, $T$ is the bulk reactor temperature, and $T_j$ is the cooling jacket temperature. The non-linear reaction rate $r_A(C_A, T)$ obeys Arrhenius kinetics:
\begin{equation}
    r_A(C_A, T) = k_0 \exp\left(-\frac{E}{R\,T}\right) C_A
    \label{eq:arrhenius}
\end{equation}
The physical state vector is $\mathbf{x} = [C_A, T, T_j]^T \in \mathbb{R}^3$, the manipulated input is the coolant volumetric flow rate $u(t) = q_j(t) \in \mathbb{R}$, and the measured process output is the reactor core temperature $y(t) = T(t)$.

\subsection{Actuator Physical Constraints and Slew-Rate Dynamics}
In physical chemical plants, control valves are constrained by finite physical stroke and actuator motor torque limits \cite{hippe2006windup, tarbouriech2009antiwindup}. The manipulated coolant flow $q_j(t)$ is bounded by hard amplitude saturation and valve travel velocity (slew-rate) limits:
\begin{align}
    u_{\min} &\le q_j(t) \le u_{\max} \label{eq:actuator_amplitude} \\
    \left|\frac{dq_j}{dt}\right| &\le S_{\max} \label{eq:actuator_slew}
\end{align}
where $u_{\min} = 0.0$~L/min represents a fully closed valve, $u_{\max} = 300.0$~L/min corresponds to full valve opening, and $S_{\max} = 80.0$~L/min$^2$ models the electro-pneumatic actuator positioning speed. Slew-rate clamping introduces an effective dynamic dead time when sudden cooling surges are commanded, severely exacerbating the risk of overshoot.

\subsection{Nominal Operating Parameters and Equilibrium}
The nominal physico-chemical parameters are compiled in Table~\ref{tab:cstr_parameters}. Under the nominal steady-state coolant flow rate $q_{j,ss} = 100.0$~L/min, the non-linear algebraic system $\mathbf{f}(\mathbf{x}_{ss}, q_{j,ss}) = \mathbf{0}$ yields the nominal operating point:
\begin{equation}
    \mathbf{x}_{ss} = \begin{bmatrix} C_{A,ss} \\ T_{ss} \\ T_{j,ss} \end{bmatrix} = \begin{bmatrix} 0.0502\text{ mol/L} \\ 396.65\text{ K} \\ 348.33\text{ K} \end{bmatrix}
    \label{eq:nominal_ss}
\end{equation}
This equilibrium corresponds to a high reactant conversion of $X_A = (C_{Af} - C_{A,ss})/C_{Af} \approx 95\%$. At this elevated conversion, the reactor operates near maximum heat release, where small positive temperature perturbations trigger exponential acceleration of the reaction rate.

\begin{table}[htbp]
\centering
\caption{Nominal physical and kinetic parameters of the simulated CSTR.}
\label{tab:cstr_parameters}
\begin{tabular}{llrl}
\toprule
Parameter & Symbol & Value & Unit \\
\midrule
Reactor volume & $V$ & $100.0$ & L \\
Cooling jacket volume & $V_j$ & $20.0$ & L \\
Process volumetric flow rate & $q$ & $100.0$ & L/min \\
Feed concentration & $C_{Af}$ & $1.0$ & mol/L \\
Feed temperature & $T_f$ & $350.0$ & K \\
Coolant feed temperature & $T_{jf}$ & $300.0$ & K \\
Pre-exponential factor & $k_0$ & $7.2 \times 10^{10}$ & min$^{-1}$ \\
Activation energy parameter & $E/R$ & $8750.0$ & K \\
Heat of reaction & $-\Delta H$ & $5.0 \times 10^4$ & cal/mol \\
Fluid heat capacity product & $\rho C_p$ & $500.0$ & cal/(L$\cdot$K) \\
Coolant heat capacity product & $\rho_j C_{pj}$ & $500.0$ & cal/(L$\cdot$K) \\
Overall heat transfer coefficient & $UA$ & $5.0 \times 10^4$ & cal/(min$\cdot$K) \\
Nominal coolant flow & $q_{j,ss}$ & $100.0$ & L/min \\
\bottomrule
\end{tabular}
\end{table}

\subsection{Thermal Controllability Boundaries and Van Heerden Analysis}
\label{subsec:thermal_controllability}
The controllability of exothermic CSTRs is fundamentally bounded by the balance between non-linear heat generation $Q_g(T)$ and dynamic heat removal $Q_r(T, q_j)$ \cite{aris1958analysis, uppal1974dynamic}:
\begin{align}
    Q_g(T) &= (-\Delta H) V r_A(C_A(T), T) \label{eq:q_gen} \\
    Q_r(T, q_j) &= q \rho C_p (T - T_f) + UA (T - T_j(q_j)) \label{eq:q_rem}
\end{align}
According to the classic Van Heerden stability criterion \cite{seborg2010process, bequette2003process}, open-loop thermal runaway occurs whenever the temperature sensitivity of heat generation exceeds that of heat removal:
\begin{equation}
    \left. \frac{\partial Q_g}{\partial T} \right|_{\mathbf{x}_{ss}} > \left. \frac{\partial Q_r}{\partial T} \right|_{\mathbf{x}_{ss}}
    \label{eq:van_heerden}
\end{equation}
At nominal steady state, the system possesses an open-loop unstable eigenvalue ($\lambda_1 = +0.284$~min$^{-1}$), requiring active feedback stabilization.

Furthermore, physical actuator saturation imposes an absolute limit on heat removal. Even at full valve capacity ($q_j = u_{\max} = 300$~L/min), the maximum achievable cooling rate is strictly bounded by:
\begin{equation}
    Q_{r, \max}(T) = q \rho C_p (T - T_f) + UA (T - T_{jf})
    \label{eq:max_cooling_limit}
\end{equation}
When the overall heat transfer capability degrades (e.g., due to polymer fouling on the heat exchange surface such that $UA < UA_{\text{crit}} \approx 0.65 UA_0$), $Q_g(T) > Q_{r, \max}(T)$ over a finite temperature interval. Under such conditions, the plant enters an unrecoverable thermal runaway regime regardless of the control law employed. This analytical boundary establishes the physical foundation for the fouling and load disturbance stress tests evaluated herein.
